% ============================================================
% 第10章 任务二：结果呈现、运营建议与管理闭环（周航正候选稿）
% 依据本地 Overleaf 工程结构整理；未自动上传。
% 与 sec09_scoring.tex 的 task2-final-rules-v1 冻结口径对齐，
% 不修改第9章的分类、公式、权重或动态政策。
% ============================================================
\section{任务二：评分结果、运营建议与管理闭环}

本章回答三个直接问题：评分能否把未来风险较高的车辆排在更需要关注的
位置，司机和管理者能否看懂分数来源，以及分数出来后应采取什么行动。
因此，结果展示与运营决策分开报告：前者使用统一折分下的未来事故与
未遂事故标签检验，后者还必须结合有效暴露、证据可信度和连续变化，
不能仅凭一次总分直接奖惩。

\subsection{冻结评分结果}

本章使用冻结版本 \texttt{task2-final-rules-v1}，覆盖 500 辆车。安全分
越高表示当前观察窗口内的行为风险越低。车队安全分均值为
\textbf{86.46}，中位数为 \textbf{88.89}，最低分为 \textbf{54.79}，
最高分为 \textbf{98.36}。其中 38 辆车因有效暴露较低而执行保守收缩，
并通过 \texttt{evidence\_flag} 单独标记，避免数据较少的车辆被误判为
优秀。

为检验界面是否易读，演示页暂以 90、80、70 分作为展示分界：90 分及
以上 165 辆，80--89.99 分 243 辆，70--79.99 分 87 辆，70 分以下
5 辆。该分段只承担视觉汇总作用，不是冻结的处罚或表扬阈值。正式运营
名单由风险分位、证据状态和动态趋势共同决定。

\begin{table}[H]
\centering
\caption{任务二冻结结果与证据边界}
\label{tab:t2-frozen-results}
\small
\begin{tabularx}{\textwidth}{lclL}
\toprule
结果项 & 数值 & 证据身份 & 管理解读 \\
\midrule
车辆数 & 500 & 冻结全量结果 & 每车均有总分、风险分位、证据状态和规则版本 \\
内部代理 AUC & 0.8158 & 未来事故/未遂标签 & 衡量排序能力，不是事故概率或预测准确率 \\
内部代理 AP & 0.2985 & 未来事故/未遂标签 & 补充反映稀有正类下的识别质量 \\
教师 Spearman & 0.552 & 任务一折外概率 & 说明排序方向的一致程度，不替代真实标签验证 \\
低暴露车辆 & 38 & 冻结证据标记 & 进入观察和补采流程，不直接处罚或参与优秀评选 \\
\bottomrule
\end{tabularx}
\end{table}

在统一内部代理协议下，规则分 AUC 为 0.8158、AP 为 0.2985，说明其
具备较好的未来风险排序能力。规则与任务一教师概率的 Spearman 相关为
0.552；教师直接风险 AUC 为 0.8652，冻结材料将规则区分能力概括为教师
参照的 94.4\%。这里的三类数字含义不同：AUC/AP回答“能否区分未来
结果”，Spearman回答“是否保留任务一排序方向”，94.4\%只表示与教师
性能的相对参照，不能替代未来标签验证。

五分位中仍存在第二组与第三组的局部倒挂，未通过最严格的逐档单调门。
这意味着总体排序有效，但中间风险段仍有重叠。运营时可以优先处理两端
的高风险与稳定优秀车辆，中间段保留人工复核，不能把相邻分档解释成
绝对等级差异。

\begin{hlbox}
\textbf{对老师最直接的结论}：评分已经能够提供有用的风险排序，但它是
管理线索而不是事故判决。低暴露、设备异常和中间分段车辆均需结合证据
复核后再采取行动。
\end{hlbox}

\subsection{结果界面与信息层级}

演示界面按“车队总览 $\to$ 司机画像 $\to$ 积分排名 $\to$ 运营动作”
组织信息，使管理者先看整体，再定位个人弱项，最后登记行动。

\begin{figure}[htbp]
  \centering
  \includegraphics[width=0.96\textwidth,height=0.76\textheight,keepaspectratio]{task2_ui_overview.png}
  \caption{车队总览、六维雷达与分数变化演示}
  \label{fig:t2-ui-overview}
\end{figure}

司机画像以冻结总分为中心，雷达图展示清醒与疲劳管理、注意力管理、
操控稳定性、安全空间管理、规则与安全配合五个行为维度；第六轴单独表示
证据可信度。证据可信度低表示判断依据不足，不等于驾驶行为更危险。
总分不再作为雷达图的一条轴，避免重复表达。变化曲线只回答近期表现
是否持续改善或恶化。

积分排名同时显示安全排名、风险分位和证据状态，防止把“安全排名第一”
与“风险排名第一”混淆。低暴露车辆保留分数供观察，但不进入正式优秀
评选。界面中总分、排名、证据状态和车队汇总来自冻结结果；由于当前
冻结交付尚未提供逐车五维拆分和 60 日动态台账，雷达、趋势及前三原因
暂用显著标注的结构演示值。正式数据接入后只替换字段，不改变页面结构。

\begin{figure}[htbp]
  \centering
  \includegraphics[width=0.96\textwidth,height=0.76\textheight,keepaspectratio]{task2_ui_ranking_operations.png}
  \caption{脱敏积分排名、四类运营名单与管理闭环演示}
  \label{fig:t2-ui-operations}
\end{figure}

\subsection{风险分层与四类运营名单}

静态分与动态状态（9.4 节）共同生成四类名单，覆盖关注、观察、改善与
表扬。每类名单都给出进入条件、处理动作和退出条件，避免名单只进不出。

\begin{table}[H]
\centering
\caption{四类运营名单及闭环规则}
\label{tab:t2-lists}
\small
\begin{tabularx}{\textwidth}{lLLL}
\toprule
名单 & 进入条件 & 处理动作 & 退出或升级条件 \\
\midrule
重点关注 & 低分且证据充分，或进入双预警区 & 24 小时内复盘前三原因，制定针对性训练 & 风险解除后降级；持续恶化则升级主管跟进 \\
观察提醒 & 中风险、低证据或短期恶化 & 补采数据、轻提醒、下一周期复核 & 证据恢复且稳定则解除；持续恶化则转重点关注 \\
改善进步 & 连续改善并解除关注状态 & 反馈具体改善行为，保留正向记录 & 持续稳定可进入表扬候选；反弹则重新观察 \\
优秀表扬 & 高分、证据充分、连续稳定且无结果门触发 & 月度表扬、经验分享、正向激励 & 证据不足或出现待核查异常则暂停资格 \\
\bottomrule
\end{tabularx}
\end{table}

两条公平性纪律贯穿四类名单。第一，低证据车辆不与正常暴露车辆竞争
正式名次，先检查设备、里程和数据覆盖，不直接处罚。第二，优秀表扬不能
由单次高分决定，必须同时满足最低有效暴露、证据可信、连续多个周期稳定
和无待核查异常；表扬规则与扣分规则同样公开。

\subsection{按维度生成可执行建议}

扣分原因直接翻译为司机话术与管理动作，五个维度分别连接到可执行的
改进路径。

\begin{table}[H]
\centering
\caption{五维度的司机建议与管理者动作}
\label{tab:t2-actions}
\small
\begin{tabularx}{\textwidth}{lLL}
\toprule
维度 & 司机建议 & 管理者动作 \\
\midrule
清醒与疲劳管理 & 合理休息，出现疲劳信号后及时停车 & 调整班次，复核夜间任务，安排疲劳教育 \\
注意力管理 & 减少电话与低头操作，驾驶前完成设备设置 & 回看具体事件，检查座舱设备交互 \\
操控稳定性 & 平稳保持车道，提前规划变道 & 结合道路和车辆条件复核，安排操控训练 \\
安全空间管理 & 增大跟车距离，降低高风险场景车速 & 检查高频路线，安排防御性驾驶培训 \\
规则与安全配合 & 遵守分场景限速，保持设备正常可用 & 复核道路限速与设备状态，处理异常申诉 \\
\bottomrule
\end{tabularx}
\end{table}

\subsection{两个可讲清楚的运营案例}

正式报告建议至少保留两个完整案例，分别覆盖风险干预和正向反馈。案例
中的车辆只使用脱敏代号，原始设备号不进入公开文档。

\paragraph{案例 A：高风险但证据充分。}
页面先显示安全总分、风险分位和证据状态，再列出前三条扣分原因。安全
经理在 24 小时内核对事件视频或轨迹场景，确认后只针对主要弱项安排
训练，并在下一周期检查分数和对应维度是否改善。若设备或场景解释不成立，
该项挂起复核，不直接转为处罚。

\paragraph{案例 B：持续改善并退出关注。}
司机连续多个周期减少主要风险行为，动态分稳定上升并退出关注区。系统
明确反馈“哪一类行为已经改善”，解除相应附加关注；若同时满足有效暴露、
证据可信和稳定性条件，可进入改善进步或优秀表扬名单。

\subsection{发布、申诉与人工可读性验收}

评分发布后按固定节奏闭环：\textbf{提醒或表扬 $\to$ 原因确认 $\to$
行动登记 $\to$ 下一周期复评 $\to$ 解除或升级}。设备异常、事件错配和
特殊任务允许挂起复核，挂起期间不把未核实证据转为正式处罚。每次动作
记录规则版本、证据来源、负责人、截止时间和复核结果，保证后续可追溯。

答辩前由不参与实现的成员独立查看页面与本章，并在 30 秒内回答：分数
方向是什么、低暴露为何不能直接评优、AUC 与教师一致性有何区别、如何
找到司机的主要弱项和下一步动作、哪些图是冻结结果以及哪些仍是结构
演示。任何一项需要作者口头补充，均视为文字或界面仍需修改。

\subsection{边界声明与版本治理}

(1) 本章指标来自内部代理标签，不等于赛事未来 40 天官方成绩；
(2) 评分用于风险排序和管理支持，风险关联不构成因果判断；
(3) 当前冻结任务二规则记录的教师版本见第 9 章（sprint-g10-final），本章按
该冻结版本报告，教师更替必须全量重跑后训练与动态层；
(4) 展示分段不自动成为奖惩阈值；
(5) 五维和趋势在冻结台账接入前只用于检验信息结构，不作为结果证据。
